% 【本文件写什么】api-v0 契约层，叶子，只写语义不写实现
% - crate 路径：os/components/wateros-ipc/ipc-waitqueue/waitqueue-api/api-v0/src/
% - 列出并说明：pub trait、核心类型、错误枚举、常量
% - 每个公开 API 的前置条件、后置条件、线程/中断上下文限制
% - 与 Linux/用户态约定的对齐点与 deliberate 差异
% - v0 版本当前行为与后续可替换 extension 点
% 【事实来源】
% - 上述 crate 下全部 .rs 源文件
% - docs/exports/public-api/ 中对应符号表，以源码为准
% 【不写什么】
% - impl 内部数据结构、汇编、具体算法步骤

\subsubsection{waitqueue-api-v0}

本 crate 定义 IPC 等待队列 trait 与类型别名，全部重导出自\code{wateros-task-api-v0}，不在 IPC 侧引入第二套等待语义。

\paragraph{类型别名。}
\code{TaskId}、\code{TaskTick}、\code{TaskWaitHandle}、\code{TaskWaitResult}、\code{WaitQueueId}与 task API 一一对应，便于 futex、pipe 与 syscall 层共享编号空间。

\paragraph{trait IpcWaitQueueOps。}
实现方须提供创建、阻塞、条件阻塞与唤醒能力；\code{wait\_current\_while}族在调度临界区内由调用方闭包复查条件，条件仍成立才让当前任务睡眠，与 Linux futex/pipe 的检查再睡模式一致。

\begin{rustcode}
pub trait IpcWaitQueueOps: Sized {
    fn new() -> Self;
    fn id(&self) -> WaitQueueId;
    fn wait_handle(&self) -> TaskWaitHandle;
    fn wait_current(&self) -> TaskWaitResult;
    fn wait_current_for_ticks(&self, timeout_ticks: TaskTick) -> TaskWaitResult;
    fn wait_current_while<F>(&self, condition: F) -> TaskWaitResult
        where F: FnOnce() -> bool;
    fn wait_current_while_for_ticks<F>(
        &self, timeout_ticks: TaskTick, condition: F) -> TaskWaitResult
        where F: FnOnce() -> bool;
    fn wake_one(&self) -> Option<TaskId>;
    fn wake_all(&self) -> usize;
}
\end{rustcode}

\paragraph{上下文限制。}
须在可阻塞的任务上下文调用\code{wait\_*}；唤醒可在持有队列引用的任意任务上下文执行。trait 未定义\code{requeue\_to}，该扩展由\code{impl-task}具体类型\code{WaitQueue}直接提供。
